\chapter*[Introduction]{Introduction}
\label{cha:introduction}
\addstarredchapter{Introduction}
\defbibheading{bibliography}[Publications]{\section*{#1}}



\section*{Contexte}


Depuis le début des années 60, les communautés de l'informatique graphique n'ont
de cesse de proposer des outils offrant toujours plus de possibilités. %
Tout d'abord avec la conception de
plans~\cite{autodeskAUTOCAD,dassaultCATIA,siemensNX,dassaultSOLIDWORKS} pour la
conception de pièces industrielles, ces outils ont gagné les domaines de la
créativité et du divertissement au travers du cinéma et des jeux
vidéos~\cite{autodesk3ds,autodeskMaya}, mais aussi l'architecture et
l'aménagement urbain~\cite{autodeskREVIT,esri2023cityengine}. %

Tous ces domaines ont un point en commun. %
Il s'agit de la construction d'objets virtuels issus d'un processus de
conception ou bien restitués à partir d'un objet réel. %


Génération automatique d'objets
\begin{itemize}
	\item construction suivant un corpus de données (génération procédurale,
	      grammaire formelle L-systems)
	\item Formalisation et cohérence des objets transformés
	\item construction manuelle d'un objet (CAO)
\end{itemize}



\section*{Problématique}



\begin{itemize}
	\item Réévaluation de modèles paramétriques
\end{itemize}



\section*{Démarche et contributions}



Intersection de trois domaines
\begin{itemize}
	\item Réévaluation : génération procédurale et CAO
	\item Génération de structures végétales : génération procédurale et
	      formalisation
	\item cohérence des objets modifiés : formalisation et CAO
\end{itemize}

Contribution
\begin{itemize}
	\item Formalisation des impacts des opérations sur un objet
	\item Formalisation et réévaluation automatique à partir de règles formelles
	\item Extension à des scripts
\end{itemize}



\section*{Organisation du manuscrit}



\begin{itemize}
	\item chapitre 1 : rappels et présentation des outils de formalisation
	\item chapitre 2 : état de l'art sur la modélisation paramétrique et la réévaluation
	\item chapitre 3 : formalisation des évènements
	\item chapitre 4 : réévaluation à partir des règles Jerboa
	\item chapitre 5 : extension aux scripts
\end{itemize}

\nocite{gaide2023automatic,gaide2023model,gaide2024reevaluation}
\printbibliography[keyword={own},title=Publications]

%%% Local Variables:
%%% mode: LaTeX
%%% TeX-master: "../../main"
%%% End:
